AI-generated software changes make it increasingly important to determine what
a change preserves, where local consistency fails to extend globally, and which
alternatives remain. This paper develops the foundations of Algebraic
Architecture Theory (AAT) from Atoms, typed primitive facts, and Laws, equations
that objects must satisfy. We make explicit the choice of what counts as
structure and which operations and laws to preserve, calling this choice a
\emph{reading}. Objects and structure-preserving morphisms are defined relative
to a reading. From finite Atom families we construct cores closed under
operations and geometries equipped with sites and coefficients.

The theory addresses five kinds of decision: gluing, diagnosis, transport,
classification of changes, and reconstruction. Its culmination is a
reconstruction theorem: the category of full geometries and all their
structure-preserving morphisms is equivalent to a category of local models
defined independently from primitive data and compatibility conditions.
Objects are recovered up to isomorphism, and morphisms between fixed endpoints
are recovered uniquely. This allows the overall structure and its changes to
be assembled from local descriptions written by separate parties.

To support these decisions, we construct a degree-one \v{C}ech obstruction to
gluing local states. Under torsor and sheaf conditions, its vanishing is
equivalent to the existence of a global state; under the conditions required
for comparison with repair semantics, it also corresponds to the existence of
a global repair. We compare diagnoses along evaluation-preserving changes of
resolution, cover, and coefficients, and give conditions for preserving first
cohomology. Invariance for all finite Law families for which both resolutions are
adequate is decidable when computable finite data are given.

Transport along exact changes is characterized by a universal property and
commutes with base change through an invertible comparison on exact pointed
pullback squares. Comparisons of two routes generated from the same square,
finite comparison diagram, and geometry factor into an invertible comparison
and an idempotent normalization. We characterize when observations alone
determine comparison preservation and classify all compatible lifts.

We encode independently defined lens (model synchronization) and protocol
semantics in this framework, proving equivalence of law satisfaction and a
one-to-one correspondence between semantics-preserving and typed morphisms.
Through these correspondences, common theorems classify and count operation-preserving
changes and yield unique extensions of morphisms from finite tables. The
corresponding Lean declarations are listed in the appendix.
